\section{Conclusion}
\label{sec:conclusion}

This paper reframed the cross-residual GNN study around biomolecular graph representation learning rather than generic benchmark breadth. The resulting picture is more precise than the earlier all-dataset narrative. Cross-residual interaction is a meaningful architectural tool, but its benefit is conditional: stronger residual reuse remains the most reliable baseline in the final topic-facing benchmark, while node-level and graph-level cross-residual variants remain competitive only on selected datasets.

The main practical outcome is therefore not that one cross-residual model dominates every benchmark. Instead, the contribution is a reusable design principle and a controlled architecture family that makes these trade-offs visible on protein-related graph tasks. In the present results, GraphResGNN is strongest on PROTEINS and ENZYMES, whereas NodeResGNN is strongest on DD. This is the right level of claim for the current evidence.

There are three immediate directions for future work. First, the topic fit can be strengthened by adding datasets that are closer to protein function prediction or to genuine multi-signal biological inference, such as protein-protein interaction datasets with ontology supervision. Second, the current graph-classification backbone can be extended toward node-level biological prediction tasks so that cross-residual mechanisms can be tested under richer supervision. Third, the architecture-level findings reported here can be paired with stronger biological features, including sequence-derived embeddings or other omics-derived signals, to move beyond structure-only benchmarks.
